\section{图论}

\subsection{最短路}

\subsubsection{Dijkstra}

\begin{lstlisting}[language=C++]
struct edge {int v, w;};

struct node {
    int dis, u;
    bool operator>(const node& a) const { return dis > a.dis; }
};

vector<edge> e[MAXN];
int dis[MAXN], vis[MAXN];
priority_queue<node, vector<node>, greater<node>> q;

void dijkstra(int n, int s) {
    memset(dis, 0x3f, (n + 1) * sizeof(int));
    memset(vis, 0, (n + 1) * sizeof(int));
    dis[s] = 0;
    q.push({0, s});
    while (!q.empty()) {
        int u = q.top().u;
        q.pop();
        if (vis[u]) continue;
        vis[u] = 1;
        for (auto ed : e[u]) {
            int v = ed.v, w = ed.w;
            if (dis[v] > dis[u] + w) {
                dis[v] = dis[u] + w;
                q.push({dis[v], v});
            }
        }
    }
}
\end{lstlisting}

\subsubsection{Bellman-Ford / SPFA}

可用于判负环

\begin{lstlisting}[language=C++]
struct edge {int v, w;};

vector<edge> e[MAXN];
int dis[MAXN], cnt[MAXN], vis[MAXN];
queue<int> q;

bool spfa(int n, int s) {
    memset(dis, 0x3f, (n + 1) * sizeof(int));
    dis[s] = 0, vis[s] = 1;
    q.push(s);
    while (!q.empty()) {
        int u = q.front();
        q.pop(), vis[u] = 0;
        for (auto ed : e[u]) {
        int v = ed.v, w = ed.w;
        if (dis[v] > dis[u] + w) {
            dis[v] = dis[u] + w;
            cnt[v] = cnt[u] + 1;  // 记录最短路经过的边数
            if (cnt[v] >= n) return false;
            // 在不经过负环的情况下，最短路至多经过 n - 1 条边
            // 因此如果经过了多于 n 条边，一定说明经过了负环
            if (!vis[v]) q.push(v), vis[v] = 1;
        }
        }
    }
    return true;
}
\end{lstlisting}

\subsubsection{Floyd}

\begin{lstlisting}[language=C++]
for (k = 1; k <= n; k++)
    for (x = 1; x <= n; x++)
        for (y = 1; y <= n; y++)
            f[x][y] = min(f[x][y], f[x][k] + f[k][y]);
\end{lstlisting}

判断任意两点之间的连通性，可用 bitset 优化至 $O(n^3/w)$

\begin{lstlisting}[language=C++]
bitset<SIZE> f[SIZE];
for (k = 1; k <= n; k++)
    for (i = 1; i <= n; i++)
        if (f[i][k]) f[i] = f[i] | f[k];
\end{lstlisting}

\subsection{基环树森林的直径}
取环+单调队列求直径. 
\begin{lstlisting}[language=C++]
// 基环树直径的两种可能：a. 不在环上 b. 环的一部分+两条链
#include<bits/stdc++.h>
using namespace std;

using ll = long long;
const int N = 2e6+10;
int n, p, tot = 1, num;
int head[N], nxt[N], ver[N], fa[N], dfn[N], s[N], q[N];
ll edge[N], ps[N], d[N], ans, res;
bool vis[N];

void add(int x, int y, int z){
    ver[++tot] = y; edge[tot] = z;
    nxt[tot] = head[x]; head[x] = tot;
}
void get_cycle(int x, int y, int z){
    ps[1] = z;
    while(y != x){
        s[++p] = y;
        ps[p+1] = edge[fa[y]];
        y = ver[fa[y]^1];
    }
    s[++p] = x;
    for(int i = 1; i <= p; i++){
        vis[s[i]] = 1;
        s[i+p] = s[i];
        ps[i+p] = ps[i];
    }
    for(int i = 1; i <= 2*p; i++)
        ps[i] += ps[i-1];
}
void dp(int x){
    vis[x] = 1;
    for(int i = head[x]; i; i = nxt[i]){
        int y = ver[i], z = edge[i];
        if(vis[y]) continue;
        dp(y);
        ans = max(ans, d[x]+d[y]+z);
        d[x] = max(d[x], d[y]+z);
    }
}
void dfs(int x){
    dfn[x] = ++num;
    for(int i = head[x]; i; i = nxt[i]){
        int y = ver[i], z = edge[i];
        if(!dfn[y]){
            fa[y] = i;
            dfs(y);
        }
        else if((i^1) != fa[x] && dfn[y] > dfn[x]){
            get_cycle(x, y, z);
        }
    }
}
int main(){
    ios::sync_with_stdio(0);
    cin >> n;
    for(int x = 1; x <= n; x++){
        int y, z; cin >> y >> z;
        add(x, y, z); add(y, x, z);
    }
    for(int i = 1; i <= n; i++) if(!dfn[i]){
        p = 0; ans = 0;
        dfs(i);
        for(int j = 1; j <= p; j++) dp(s[j]);
        int l = 1, r = 0;
        for(int j = 1; j <= 2*p; j++){
            while(l <= r && j-q[l] >= p) l++;
            if(l <= r) ans = max(ans, d[s[j]]+d[s[q[l]]]+ps[j]-ps[q[l]]);
            while(l <= r && d[s[q[r]]]-ps[q[r]] <= d[s[j]]-ps[j]) r--;
            q[++r] = j; 
        }
        res += ans;
    }
    cout << res << endl;
    return 0;
}
\end{lstlisting}

\subsection{树链剖分}
轻重链剖分
\begin{lstlisting}[language=C++]
struct HLD {
    int n;
    vector<vector<int>> adj;
    vector<int> dep, siz, parent, in, out, ord, top;
    int cur;

    HLD(int n) {
        init(n);
    }
    void init(int n) {
        this->n = n;
        adj.assign(n, {});
        dep.resize(n);
        siz.resize(n);
        parent.resize(n);
        in.resize(n);
        out.resize(n);
        ord.resize(n);
        top.resize(n);
        cur = 0;
    }
    void addEdge(int u, int v) {
        adj[u].push_back(v);
        adj[v].push_back(u);
    }
    void work(int root) {
        top[root] = root;
        dep[root] = 1;
        parent[root] = 0;
        dfs1(root, 0);
        dfs2(root, 0);
    }
    void dfs1(int u, int pa) {
        siz[u] = 1;
        for (auto &v : adj[u]) {
            if (v == pa) {
                continue;
            }
            parent[v] = u;
            dep[v] = dep[u] + 1;
            dfs1(v, u);
            siz[u] += siz[v];
            if (siz[v] > siz[adj[u][0]]) {
                swap(v, adj[u][0]);
            }
        }
    }
    void dfs2(int u, int pa) {
        in[u] = ++cur;
        ord[in[u]] = u;
        for (auto v : adj[u]) {
            if (v == pa) {
                continue;
            }
            top[v] = (v == adj[u][0] ? top[u] : v);
            dfs2(v, u);
        }
        out[u] = cur;
    }
    int lca(int u, int v) {
        while (top[u] != top[v]) {
            if (dep[top[u]] > dep[top[v]]) {
                swap(u, v);
            }
            v = parent[top[v]];
        }
        return dep[u] > dep[v] ? v : u;
    }
};

// 路径修改
auto pathModify = [&](int u, int v, int k) {
    while (t.top[u] != t.top[v]) {
        if (t.dep[t.top[u]] > t.dep[t.top[v]]) {
            swap(u, v);
        }
        seg.rangeModify(t.in[t.top[v]], t.in[v], {k});
        v = t.parent[t.top[v]];
    }
    if (t.dep[u] > t.dep[v]) {
        swap(u, v);
    }
    seg.rangeModify(t.in[u], t.in[v], {k});
};

// 路径查询
auto pathQuery = [&](int u, int v) {
    int res = 0;
    while (t.top[u] != t.top[v]) {
        if (t.dep[t.top[u]] > t.dep[t.top[v]]) {
            swap(u, v);
        }
        res += seg.rangeQuery(t.in[t.top[v]], t.in[v]).sum, res %= mod;
        v = t.parent[t.top[v]];
    }
    if (t.dep[u] > t.dep[v]) {
        swap(u, v);
    }
    res += seg.rangeQuery(t.in[u], t.in[v]).sum, res %= mod;
    return res;
};

// 子树修改
auto subtreeModify = [&](int u, int k) {
    seg.rangeModify(t.in[u], t.out[u], {k});
};

// 子树查询
auto subtreeQuery = [&](int u) {
    return seg.rangeQuery(t.in[u], t.out[u]).sum % mod;
};
// by lyh
\end{lstlisting}

\subsection{LCA}

\subsection{无向图连通性}

\subsubsection{割点和桥}

割点判定法则：
\[dfn[x]\le low[y]\]
其中 $low[y]$ 表示 \underline{子树}以及 \underline{通过一条非搜索树上的边能够到达的子树外的点} 的最小 $dfn$。特别地，对于根节点，要有两个点满足上述条件。由于是小于等于，父亲节点实际上可取。
\begin{lstlisting}[language=C++]
void Tarjan(int u, int pa) {
    low[u] = dfn[u] = ++dfn_cnt;
    for (auto v : adj[u]) {
        if (v == pa) continue;
        if (!dfn[v]) {
            if (pa == 0) {
                son_cnt++; // 若u是搜索起点, 更新儿子数量
            }
            Tarjan(v, u);
            low[u] = min(low[u], low[v]);

            if (pa != 0 && low[v] >= dfn[u] && !is_cut[u]) { 
                cut_cnt++;
                is_cut[u] = true;
            }
        } else {
            low[u] = min(low[u], dfn[v]); 
        }
    }
    if (pa == 0 && son_cnt >= 2 && !is_cut[u]) { // 搜索起点有两个及以上的儿子
        cut_cnt++;
        is_cut[u] = true;
    }
}
//by lyh
\end{lstlisting}

\begin{lstlisting}[language=C++]
void tarjan(int x){
    dfn[x] = low[x] = ++tot;
    int cnt = 0;
    for(auto y: e[x]) if(!dfn[y]){
        tarjan(y);
        low[x] = min(low[x], low[y]);
        if(dfn[x] <= low[y]){
            ++cnt;
            if(x != root || cnt > 1) cut[x] = 1;
        }
    }
    else low[x] = min(low[x], dfn[y]);
}
// if(!dfn[x]) root = x, tarjan(x);
// by zzy
\end{lstlisting}

割边判定法则：
\[dfn[x]< low[y]\]
父亲节点不可取，注意特判有重边的情况。
\begin{lstlisting}[language=C++]
void Tarjan(int u, int lst) {
    low[u] = dfn[u] = ++dfn_cnt;
    for (auto [v, id] : adj[u]) {
        if (id == lst) {
            continue;
        }
        if (!dfn[v]) {
            Tarjan(v, id);
            low[u] = min(low[u], low[v]);
            
            if (low[v] > dfn[u]) {
                bridge_cnt++;
                is_bridge[id] = true;
            }
        } else {
            low[u] = min(low[u], dfn[v]);
        }
    }
}
// by lyh
\end{lstlisting}

\subsubsection{e-DCC 和 v-DCC}

边双连通分量
\begin{lstlisting}[language=C++]
struct EBCC {
    int n;
    vector<vector<array<int, 2>>> adj;
    vector<int> stk;
    vector<int> dfn, low, bel;
    int cur, cnt;

    EBCC(int n) {
        init(n);
    }
    void init(int n) {
        this->n = n;
        adj.assign(n, {});
        dfn.assign(n, 0);
        low.resize(n);
        bel.assign(n, 0);
        stk.clear();
        cur = cnt = 0;
    }
    void addEdge(int u, int v, int id) {
        adj[u].push_back({v, id});
        adj[v].push_back({u, id});
    }
    void Tarjan(int u, int lst) {
        dfn[u] = low[u] = ++cur;
        stk.push_back(u);

        for (auto [v, id] : adj[u]) {
            if (id == lst) {
                continue;
            }
            if (!dfn[v]) {
                Tarjan(v, id);
                low[u] = min(low[u], low[v]);
            } else {
                low[u] = min(low[u], dfn[v]);
            }
        }

        if (low[u] == dfn[u]) {
            cnt++;
            int v;
            do {
                v = stk.back();
                stk.pop_back();
                bel[v] = cnt;
            } while (v != u);
        }
    }
    vector<int> work() {
        for (int i = 1; i < n; i++) {
            if (!dfn[i]) {
                Tarjan(i, 0);
            }
        }
        return bel;
    }
};
\end{lstlisting}

点双连通分量
\begin{lstlisting}[language=C++]
struct VBCC {
    int n;
    vector<vector<int>> adj;
    vector<int> stk;
    vector<int> dfn, low;
    int cur, cnt;

    VBCC(int n) {
        init(n);
    }
    void init(int n) {
        this->n = n;
        adj.assign(n, {});
        dfn.assign(n, 0);
        low.resize(n);
        cur = cnt = 0;
    }
    void addEdge(int u, int v) {
        adj[u].push_back(v);
        adj[v].push_back(u);
    }
    void Tarjan(int u, int pa) {
        dfn[u] = low[u] = ++cur;
        stk.push_back(u);
        
        for (auto v : adj[u]) {
            if (v == pa) {
                continue;
            }
            if (!dfn[v]) {
                Tarjan(v, u);
                low[u] = min(low[u], low[v]);

                if (low[v] >= dfn[u]) {
                    vbcc.emplace_back();
                    cnt++;
                    int y;
                    do {
                        y = stk.back();
                        stk.pop_back();
                        vbcc[cnt].push_back(y);
                    } while (y != v);
                    vbcc[cnt].push_back(u);
                }
            } else {
                low[u] = min(low[u], dfn[v]);
            }
        }

        if (pa == 0 && adj[u].size() == 0) {
            vbcc.emplace_back();
            cnt++;
            vbcc[cnt].push_back(u);
        }
    }
    auto work() {
        for (int i = 1; i < n; i++) {
            if (!dfn[i]) {
                Tarjan(i, 0);
            }
        }
        return vbcc;
    }
};
// by lyh
\end{lstlisting}

\subsection{有向图连通性}

Tarjan求强连通分量
\begin{lstlisting}[language=C++]
struct SCC {
    int n;
    vector<vector<int>> adj;
    vector<int> stk;
    vector<int> dfn, low, bel;
    int cur, cnt;
    
    SCC(int n) {
        init(n);
    }
    void init(int n) {
        this->n = n;
        adj.assign(n, {});
        dfn.assign(n, 0);
        low.resize(n);
        bel.assign(n, 0);
        stk.clear();
        cur = cnt = 0;
    }
    void addEdge(int u, int v) {
        adj[u].push_back(v);
    }
    void Tarjan(int u) {
        dfn[u] = low[u] = ++cur;
        stk.push_back(u);
        
        for (auto v : adj[u]) {
            if (!dfn[v]) {
                Tarjan(v);
                low[u] = min(low[u], low[v]);
            } else if (!bel[v]) {
                low[u] = min(low[u], dfn[v]);
            }
        }
        
        if (dfn[u] == low[u]) {
            cnt++;
            int v;
            do {
                v = stk.back();
                stk.pop_back();
                bel[v] = cnt;
            } while (v != u);
        }
    }
    vector<int> work() {
        for (int i = 1; i < n; i++) {
            if (!dfn[i]) {
                Tarjan(i);
            }
        }
        return bel;
    }
};
// by lyh
\end{lstlisting}

\subsection{欧拉回路}
有向图版本
\begin{lstlisting}[language=C++]
void solve() {
    int n, m;
    cin >> n >> m;
    vector<vector<int>> e(n + 1);
    vector<int> in(n + 1), out(n + 1);
    for (int i = 1; i <= m; i++) {
        int u, v;
        cin >> u >> v;
        e[u].push_back(v);
        out[u]++;
        in[v]++;
    }
    
    int cnt1 = 0, cnt2 = 0, s = 1;
    for (int i = 1; i <= n; i++) {
        sort(e[i].begin(), e[i].end());
        if (out[i] != in[i]) {
            if (out[i] == in[i] + 1) {
                cnt1++;
                s = i;
            } else if (out[i] + 1 == in[i]) {
                cnt2++;
            } else {
                cout << "No\n";
                return;
            }
        }
    }
    if (!(cnt1 == 1 && cnt2 == 1 || cnt1 == 0 && cnt2 == 0)) {
        cout << "No\n";
        return;
    }

    vector<int> cur(n + 1);
    stack<int> st;
    auto dfs = [&](auto self, int u) -> void {
        for (int i = cur[u]; i < e[u].size(); i = cur[u]) {
            cur[u] = i + 1;
            int v = e[u][i];
            self(self, v);
        }
        st.push(u);
    };
    dfs(dfs, s);

    while (st.size()) {
        cout << st.top() << " ";
        st.pop();
    }
}
\end{lstlisting}
无向图版本
\begin{lstlisting}[language=C++]
void solve() {
    int n = 500, m;
    cin >> m;
    vector<vector<array<int, 2>>> e(n + 1);
    vector<int> d(n + 1);
    int s = n;
    for (int i = 1; i <= m; i++) {
        int u, v;
        cin >> u >> v;
        e[u].push_back({v, i});
        e[v].push_back({u, i});
        d[u]++;
        d[v]++;
        s = min(s, min(u, v));
    }
    
    for (int i = n; i >= 1; i--) {
        sort(e[i].begin(), e[i].end());
        if (d[i] % 2 == 1) {
            s = i;
        }
    }

    vector<bool> vis(m + 1);
    vector<int> cur(n + 1);
    stack<int> st;
    auto dfs = [&](auto &&self, int u) -> void {
        for (int &j = cur[u]; j < e[u].size(); j++) {
            auto [v, i] = e[u][j];
            if (!vis[i]) {
                vis[i] = true;
                self(self, v);
                // ans.push_back(i);
            }
        }
        st.push(u);
    };
    dfs(dfs, s);

    while (st.size()) {
        cout << st.top() << "\n";
        st.pop();
    }
}
\end{lstlisting}
\subsection{二分图}

\subsection{网络流}

\subsubsection{最大流}

\begin{lstlisting}[language=C++]
template<class T>
struct MaxFlow {
    struct _Edge {
        int to;
        T cap;
    };
    int n;
    vector<_Edge> e;
    vector<vector<int>> g;
    vector<int> cur, h;
    MaxFlow(int n) {
        init(n);
    }
    void init(int n) {
        this->n = n;
        e.clear();
        g.assign(n, {});
        cur.resize(n);
        h.resize(n);
    }
    bool bfs(int s, int t) {
        h.assign(n, -1);
        queue<int> q;
        h[s] = 0;
        q.push(s);
        while (!q.empty()) {
            auto u = q.front();
            q.pop();
            for (auto i : g[u]) {
                auto [v, c] = e[i];
                if (c > 0 && h[v] == -1) {
                    h[v] = h[u] + 1;
                    if (v == t) {
                        return true;
                    }
                    q.push(v);
                }
            }
        }
        return false;
    }
    T dfs(int u, int t, T f) {
        if (u == t) {
            return f;
        }
        auto r = f;
        for (int &i = cur[u]; i < int(g[u].size()); i++) {
            auto j = g[u][i];
            auto [v, c] = e[j];
            if (c > 0 && h[v] == h[u] + 1) {
                auto a = dfs(v, t, min(r, c));
                e[j].cap -= a;
                e[j ^ 1].cap += a;
                r -= a;
                if (r == 0) {
                    return f;
                }
            }
        }
        return f - r;
    }
    void addEdge(int u, int v, T cap) {
        g[u].push_back(e.size());
        e.push_back({v, cap});
        g[v].push_back(e.size());
        e.push_back({u, 0});
    }
    T flow(int s, int t) {
        T ans = 0;
        while (bfs(s, t)) {
            cur.assign(n, 0);
            ans += dfs(s, t, inf);
        }
        return ans;
    }
    vector<bool> minCut() {
        vector<bool> c(n);
        for (int i = 0; i < n; i++) {
            c[i] = (h[i] != -1);
        }
        return c;
    }
    struct Edge {
        int from;
        int to;
        T cap;
        T flow;
    };
    vector<Edge> edges() {
        vector<Edge> res;
        for (int i = 0; i < e.size(); i += 2) {
            res.push_back({e[i + 1].to, e[i].to, e[i].cap + e[i + 1].cap, e[i + 1].cap});
        }
        return res;
    }
};
// by lyh
\end{lstlisting}

\subsubsection{最小费用最大流}

\begin{lstlisting}[language=C++]
template<class T>
struct MinCostFlow {
    struct _Edge {
        int to;
        T cap;
        T cost;
    };
    int n;
    vector<_Edge> e;
    vector<vector<int>> g;
    vector<T> h, dis;
    vector<int> pre;
    bool Dijkstra(int s, int t) {
        dis.assign(n, inf);
        pre.assign(n, -1);
        priority_queue<pair<T, int>, vector<pair<T, int>>, greater<pair<T, int>>> q;
        dis[s] = 0;
        q.push({0, s});
        while (!q.empty()) {
            auto [d, u] = q.top();
            q.pop();
            if (dis[u] < d) {
                continue;
            }
            for (auto id : g[u]) {
                auto [v, cap, cost] = e[id];
                if (cap > 0 && dis[v] > d + h[u] - h[v] + cost) {
                    dis[v] = d + h[u] - h[v] + cost;
                    pre[v] = id;
                    q.push({dis[v], v});
                }
            }
        }
        return dis[t] != inf;
    }
    MinCostFlow() {}
    MinCostFlow(int n_) {
        init(n_);
    }
    void init(int n_) {
        n = n_;
        e.clear();
        g.assign(n, {});
    }
    void addEdge(int u, int v, T cap, T cost) {
        g[u].push_back(e.size());
        e.push_back({v, cap, cost});
        g[v].push_back(e.size());
        e.push_back({u, 0, -cost});
    }
    pair<T, T> flow(int s, int t) {
        T flow = 0, cost = 0;
        h.assign(n, 0);
        while (Dijkstra(s, t)) {
            for (int i = 0; i < n; i++) {
                h[i] += dis[i];
            }
            T aug = inf;
            for (int i = t; i != s; i = e[pre[i] ^ 1].to) {
                aug = min(aug, e[pre[i]].cap);
            }
            for (int i = t; i != s; i = e[pre[i] ^ 1].to) {
                e[pre[i]].cap -= aug;
                e[pre[i] ^ 1].cap += aug;
            }
            flow += aug;
            cost += aug * h[t];
        }
        return {flow, cost};
    }
    struct Edge {
        int from;
        int to;
        T cap;
        T cost;
        T flow;
    };
    vector<Edge> edges() {
        vector<Edge> res;
        for (int i = 0; i < e.size(); i += 2) {
            Edge ed = {e[i + 1].to, e[i].to, e[i].cap + e[i + 1].cap, e[i].cost, e[i + 1].cap};
            res.push_back(ed);
        }
        return res;
    }
};
// by lyh
\end{lstlisting}
